\documentclass[11pt]{article}
\usepackage[margin=1in]{geometry}
\usepackage{amsmath,amssymb,amsthm}
\usepackage[T1]{fontenc}
\usepackage{lmodern}
\usepackage{graphicx}
\usepackage[hidelinks]{hyperref}
\usepackage{xcolor}

\newtheorem{theorem}{Theorem}
\newtheorem{lemma}[theorem]{Lemma}
\newtheorem{remark}[theorem]{Remark}
\newcommand{\R}{\mathbb R}
\newcommand{\N}{\mathbb N}
\newcommand{\F}{\mathcal F}
\newcommand{\spec}{\operatorname{spec}}

\title{Path-Connectedness of Scale-Invariant\\
Floer-Hessian Growth Strata}
\author{Candidate partial result for arXiv:2408.00269}
\date{11 August 2026}

\begin{document}
\maketitle

\begin{center}
\fcolorbox{black}{gray!8}{\parbox{0.88\linewidth}{
\textbf{Status: candidate partial result, likely valid.}
Question 1.7 remains open for general shift-invariant, non-scale-invariant
pair growth.}}
\end{center}

\section{Source question and notation}

Frauenfelder and Weber ask the following in \cite[Question 1.7]{FW}, on
arXiv PDF page 6.

\begin{quote}
\textbf{Question 1.7.} Given shift-invariant growth types
$\mathfrak a,\mathfrak b\in\mathfrak G_1$, is the space
$\F_h^{\mathfrak a\mathfrak b}$ connected?
\end{quote}

\begin{figure}[h]
  \centering
  \includegraphics[width=\textwidth]{figures/open_problem_crop.png}
  \caption{Question 1.7 and its context in arXiv:2408.00269, PDF page 6.}
\end{figure}

We use the definitions of \cite{FW}.  Thus
\[
 H_r=\ell^2_{h^r},\qquad r\in\R,
\]
and a weak Hessian is an $H_0$-symmetric Fredholm operator of index zero
$A:H_1\to H_0$ whose restriction $A:H_2\to H_1$ is also Fredholm of
index zero.  The subspace $\F_h^{\mathfrak a\mathfrak b}$ consists of the
Floer-type weak Hessians whose negative and positive eigenvalue growth types
are $\mathfrak a$ and $\mathfrak b$, in the translation-invariant sense of
\cite[Section 1.1]{FW}.

A growth type $[h]$ is \emph{scale-invariant} when it has a representative
satisfying
\begin{equation}\label{eq:doubling}
 h(2n)\le C h(n)\qquad(n\in\N)
\end{equation}
for some $C$.  This is stronger than shift invariance; polynomial growth is
scale-invariant, whereas exponential growth is shift-invariant but not
scale-invariant \cite[Section 2.2.2]{FW}.

\section{Proof intuition}

Translate two Hessians slightly so that both become invertible.  Their
positive eigenvalues have the same growth type, as do their negative
eigenvalues.  Match positive eigenvectors to positive eigenvectors by rank,
and do the same on the negative side.  The resulting orthogonal operator is
not merely an $H_0$-unitary: the graph norms of the two Hessians show that it
and its inverse are bounded on $H_1$ and $H_2$.

The key use of scale invariance is topological.  Applied to the Hilbert pair
$(H_0,H_2)$, whose weight is $h^2$, the mild-spectrum Kuiper theorem connects
the matching operator to the identity through orthogonal maps bounded on
$H_2$.  Conjugating along this path aligns the two eigenbases while staying
inside the weak-Hessian topology.  In the common eigenbasis, linearly
interpolate corresponding eigenvalues.  Since matched eigenvalues have the
same sign and comparable magnitudes, no eigenvalue reaches zero and both
Fredholm-level estimates remain uniform.

\section{Two lemmas}

\begin{lemma}[Scale-orthogonal path]\label{lem:kuiper}
Assume \eqref{eq:doubling}.  If $U$ is orthogonal on $H_0$ and $U,U^{-1}$
restrict to bounded operators on $H_2$, then there is a path $(U_t)_{0\le t\le
1}$ from the identity to $U$ such that every $U_t$ is orthogonal on $H_0$,
$U_t,U_t^{-1}$ are bounded on $H_2$, and the path is continuous in the pair
operator topology.  Consequently it is continuous on $H_1$ as well.
\end{lemma}

\begin{proof}
The Hilbert pair $(H_0,H_2)$ is isometric to
$(\ell^2,\ell^2_{h^2})$.  Its pair weight is mild because
\[
 h(2n)^2\le C^2h(n)^2.
\]
The mild-spectrum Kuiper theorem
\cite[Proposition 33.1.5]{KM} says that the group of pair-unitaries is
contractible.  Its proof is Kuiper's argument in the pair-bounded operator
algebra; the same proof over the real field replaces unitary maps by
orthogonal maps, as in the real form of Kuiper's theorem \cite{Kuiper}.
This supplies $(U_t)$.  Boundedness and norm-continuity on $H_1$, the midpoint
space of the Hilbert scale, follow by interpolation between $H_0$ and $H_2$.
\end{proof}

\begin{lemma}[Sign-preserving eigenbasis matching]\label{lem:matching}
Let $A,B\in\F_h^{\mathfrak a\mathfrak b}$ be invertible.  There is an
$H_0$-orthogonal operator $U$ which maps the positive eigenvectors of $A$ to
those of $B$ in increasing spectral order and likewise on the negative side,
and for which $U,U^{-1}$ are bounded on $H_1$ and $H_2$.
\end{lemma}

\begin{proof}
Choose $H_0$-orthonormal eigenbases
\[
 Av_n^+=\alpha_n^+v_n^+,quad Av_n^-=-\alpha_n^-v_n^-,
 \qquad
 Bw_n^+=\beta_n^+w_n^+,quad Bw_n^-=-\beta_n^-w_n^-,
\]
where all displayed eigenvalue magnitudes are positive and are nondecreasing,
with multiplicities.  Define $Uv_n^\pm=w_n^\pm$.  This is orthogonal on
$H_0$.

Equality of the signed growth types gives a constant $K\ge1$ such that
\begin{equation}\label{eq:compare}
 K^{-1}\beta_n^\pm\le\alpha_n^\pm\le K\beta_n^\pm
 \qquad(n\in\N).
\end{equation}
If the source convention records squared eigenvalues, take square roots of
its comparison constants to obtain \eqref{eq:compare}.

Because $A:H_1\to H_0$ is an isomorphism, its graph norm is equivalent to
the fixed $H_1$ norm.  Hence, for finite eigenvector sums,
\begin{equation}\label{eq:level1A}
 \|x\|_1^2\asymp\|Ax\|_0^2
 =\sum_{n,\varepsilon\in\{+,-\}}
 (\alpha_n^\varepsilon)^2|x_n^\varepsilon|^2.
\end{equation}
The analogous formula holds for $B$ in the $w$-basis.  Equations
\eqref{eq:compare}--\eqref{eq:level1A} show that $U$ and $U^{-1}$ are
bounded on $H_1$.

The regularity theorem of \cite{FW} says that invertibility on
$H_1\to H_0$ is equivalent to invertibility of the restriction
$A:H_2\to H_1$.  Applying the level-one graph-norm equivalence twice gives
\begin{equation}\label{eq:level2A}
 \|x\|_2\asymp\|Ax\|_1\asymp\|A^2x\|_0,
\end{equation}
and therefore, again first on finite sums,
\[
 \|x\|_2^2\asymp
 \sum_{n,\varepsilon\in\{+,-\}}
 (\alpha_n^\varepsilon)^4|x_n^\varepsilon|^2.
\]
The corresponding formula for $B$, together with \eqref{eq:compare}, proves
boundedness of $U$ and $U^{-1}$ on $H_2$.  Density extends the estimates to
the full spaces.
\end{proof}

\section{Partial answer to Question 1.7}

\begin{theorem}\label{thm:main}
Let the pair growth type $[h]$ be scale-invariant, and let
$\mathfrak a,\mathfrak b$ be shift-invariant growth types.  Every nonempty
space $\F_h^{\mathfrak a\mathfrak b}$ is path-connected in the topology
induced by
\[
 \mathcal L(H_1,H_0)\cap\mathcal L(H_2,H_1).
\]
\end{theorem}

\begin{proof}
Take $A,B\in\F_h^{\mathfrak a\mathfrak b}$.  The spectrum of each weak
Hessian is discrete and has nonempty resolvent.  Choose
$\lambda_A\notin\spec A$ and $\lambda_B\notin\spec B$.  The translation
paths
\[
 s\longmapsto A-s\lambda_A\iota,
 \qquad
 s\longmapsto B-s\lambda_B\iota
\]
remain in $\F_h^{\mathfrak a\mathfrak b}$ by the translation invariance
proved in \cite[Section 1.1]{FW}, and their endpoints
$A'=A-\lambda_A\iota$ and $B'=B-\lambda_B\iota$ are invertible.  It is
therefore enough to connect $A'$ and $B'$.

Apply Lemma~\ref{lem:matching} and let $U$ be the sign-preserving eigenbasis
matching operator.  Lemma~\ref{lem:kuiper} supplies a path $U_t$ from the
identity to $U$.  Put
\[
 A_t=U_tA'U_t^{-1}.
\]
Every $A_t$ is $H_0$-symmetric because $U_t$ is $H_0$-orthogonal.  Since
$U_t$ and its inverse preserve $H_1$ and $H_2$, conjugation preserves
Fredholm invertibility on both $H_1\to H_0$ and $H_2\to H_1$.  The signed
spectrum is unchanged, so $A_t\in\F_h^{\mathfrak a\mathfrak b}$.  The
operator path is continuous in both defining norms.  Thus $A'$ is connected
to $\widetilde A:=UA'U^{-1}$.

By construction, $\widetilde A$ and $B'$ are diagonal in the same
$H_0$-orthonormal eigenbasis, with corresponding eigenvalues of the same
sign.  Write these signed eigenvalues as $a_j$ and $b_j$.  Their magnitudes
obey
\[
 K^{-1}|b_j|\le |a_j|\le K|b_j|.
\]
For $0\le s\le1$, define
\[
 C_s=(1-s)\widetilde A+sB'.
\]
On the common eigenbasis the eigenvalues are
$c_j(s)=(1-s)a_j+sb_j$.  Same signs and the comparison above imply
\[
 \min\{1,K^{-1}\}|b_j|
 \le |c_j(s)|
 \le \max\{1,K\}|b_j|.
\]
Consequently $C_s$ is an isomorphism $H_1\to H_0$ and, using the squared
graph-norm comparison, an isomorphism $H_2\to H_1$.  It is
$H_0$-symmetric, of Floer type, and its positive and negative eigenvalue
growth types remain $\mathfrak a$ and $\mathfrak b$.  Hence
$C_s\in\F_h^{\mathfrak a\mathfrak b}$ throughout.  This path joins
$\widetilde A$ to $B'$.

Concatenate these three paths with the reverse translation from $B'$ to $B$.
The resulting path lies entirely in
$\F_h^{\mathfrak a\mathfrak b}$.
\end{proof}

\section{Limitations and novelty check}

The theorem does not prove connectivity for all shift-invariant growth.  The
scale-Kuiper input requires mild spectrum.  For $(H_0,H_2)$ that hypothesis is
precisely supplied by scale invariance of $h$; shift invariance alone does not
imply it.  In particular the exponential example $h(n)=e^n$ from \cite{FW}
lies outside the proof.

A bounded web and arXiv search on 11 August 2026 used the exact paper title,
the wording of Question 1.7, and the phrases ``Floer Hessian connected'',
``scale invariant Floer Hessian'', and ``fixed signed growth stratum''.  It
located the source paper and later citations concerned with different
Fredholm questions, but no paper claiming either a solution to Question 1.7
or the scale-invariant theorem above.  Novelty confidence is therefore
bounded rather than definitive.

\section{Verification and human-review focus}

The proof was checked at four sensitive points: mildness of the pair
$(H_0,H_2)$; the $H_2$ graph-norm identity; real-field applicability of the
scale-Kuiper argument; and preservation of both weak-Hessian levels during
eigenvalue interpolation.  No computational experiment is used.

Human review should focus first on Lemma~\ref{lem:kuiper}, specifically the
real orthogonal form of the cited pair-Kuiper theorem, and then on
\eqref{eq:level2A}.  Once those are accepted, the path construction is
explicit.

\begin{thebibliography}{9}
\bibitem{FW}
U. Frauenfelder and J. Weber,
\emph{Growth of eigenvalues of Floer Hessians},
arXiv:2408.00269, 2024.

\bibitem{KM}
P. B. Kronheimer and T. S. Mrowka,
\emph{Monopoles and Three-Manifolds},
Cambridge University Press, 2007; Proposition 33.1.5.

\bibitem{Kuiper}
N. H. Kuiper,
\emph{The homotopy type of the unitary group of Hilbert space},
Topology 3 (1965), 19--30.
\end{thebibliography}

\end{document}
